\documentclass[12pt]{article}
\usepackage[utf8]{inputenc}
\usepackage[russian]{babel}
\usepackage{amsmath}
\usepackage{amssymb}
\usepackage{geometry}
\geometry{a4paper, margin=2cm}

\begin{document}

\section*{Задача 1.9}

\textbf{Условие:} Построить распознаватель языка
$L=\{a^k b^k c^k:k\geq 0\}$. Если слово принадлежит языку,
лента должна очиститься; иначе машина должна перейти в состояние
\texttt{ERROR}.

\section*{Программа}

\begin{verbatim}
alphabet = [#, a, b, c, A, B, C]
tapes = [ right ]

START:
    # -> #R FIND_A
    a -> aH ERROR
    b -> bH ERROR
    c -> cH ERROR

FIND_A:
    A -> AR FIND_A
    a -> AR FIND_B
    B -> BR CHECK_REST
    b -> bR CHECK_REST
    C -> CR CHECK_REST
    c -> cR CHECK_REST
    # -> #H STOP

FIND_B:
    a -> aR FIND_B
    B -> BR FIND_B
    b -> BR FIND_C
    A -> AH ERROR
    C -> CH ERROR
    c -> cH ERROR
    # -> #H ERROR

FIND_C:
    b -> bR FIND_C
    B -> BR FIND_C
    C -> CR FIND_C_AFTER_C
    c -> CL RETURN_LEFT
    a -> aH ERROR
    A -> AH ERROR
    # -> #H ERROR

FIND_C_AFTER_C:
    C -> CR FIND_C_AFTER_C
    c -> CL RETURN_LEFT
    a -> aH ERROR
    A -> AH ERROR
    b -> bH ERROR
    B -> BH ERROR
    # -> #H ERROR

RETURN_LEFT:
    a -> aL RETURN_LEFT
    b -> bL RETURN_LEFT
    c -> cL RETURN_LEFT
    A -> AL RETURN_LEFT
    B -> BL RETURN_LEFT
    C -> CL RETURN_LEFT
    # -> #R FIND_A

CHECK_REST:
    A -> AR CHECK_REST
    B -> BR CHECK_REST
    C -> CR CHECK_REST
    a -> aH ERROR
    b -> bH ERROR
    c -> cH ERROR
    # -> #L CLEANUP

CLEANUP:
    A -> #L CLEANUP
    B -> #L CLEANUP
    C -> #L CLEANUP
    a -> aH ERROR
    b -> bH ERROR
    c -> cH ERROR
    # -> #H STOP
\end{verbatim}

\subsection*{Точная сложность в худшем случае}

Пусть $n=|w|$ --- длина слова без граничных символов;
каждый переход, включая переход в \texttt{STOP} или \texttt{ERROR},
считается за один шаг. Для пустого слова
\[
T(\varepsilon)=2.
\]

Рассмотрим слово, содержащие начальные блоки
$w=a^p b^q c^r v$, где $p,q,r\geq 1$ и $v=\varepsilon$ либо
первый символ $v$ равен $a$ или $b$. Пусть
$k=\min(p,q,r)$. За $j$-ю успешную итерацию машина
просматривает $j$ символов в $\texttt{FIND\_A}$,
$p$ символов в $\texttt{FIND\_B}$,
$q$ символов в $\texttt{FIND\_C}$ и
$p+q+j$ символов при возвращении к левой границе. Поэтому
\[
T_j=j+p+q+(p+q+j)=2(p+q+j).
\]
Суммарная стоимость $k$ успешных итераций равна
\[
R(p,q,k)=\sum_{j=1}^{k}2(p+q+j)
       =2k(p+q)+k(k+1).
\]

После $k$ итераций машина либо проверяет остаток, либо пытается
выполнить ещё одну итерацию. Точная стоимость этой заключительной
части обозначается $U$. Если $v\ne\varepsilon$, то
\[
U_{\mathrm{off}}(p,q,r)=
\begin{cases}
2k+1, & p=k,\ q>k,\\
p+k+1, & p>k,\ q=k,\\
3k+1, & p=q=k,\ r>k,\\
p+q+k+1, & p>k,\ q>k,\ r=k,\\
3k+1, & p=q=r=k.
\end{cases}
\]
Здесь $v\ne\varepsilon$ означает, что после блока $c^r$ остался
хотя бы один символ. Если вход заканчивается ровно после блока
$c^r$, то заключительная часть имеет ту же стоимость, кроме
случая $p=q=r=k$: в этом случае машина очищает все $3k$ символов,
и
\[
U_{\mathrm{end}}(p,q,r)=
\begin{cases}
2k+1, & p=k,\ q>k,\\
p+k+1, & p>k,\ q=k,\\
3k+1, & p=q=k,\ r>k,\\
p+q+k+1, & p>k,\ q>k,\ r=k,\\
6k+2, & p=q=r=k.
\end{cases}
\]
Таким образом,
\[
T_{\mathrm{off}}(p,q,r)=1+R(p,q,k)+U_{\mathrm{off}}(p,q,r),
\]
\[
T_{\mathrm{end}}(p,q,r)=1+R(p,q,k)+U_{\mathrm{end}}(p,q,r).
\]
Единица в формулах учитывает переход из \texttt{START}.

Каждая успешная итерация использует по одному символу каждого типа,
поэтому число полных итераций не превосходит $\lfloor n/3\rfloor$.
Покажем, что меньшее число итераций не даёт максимума. Если выполнено
$t<\lfloor n/3\rfloor$ полных итераций, то заключительная часть имеет
стоимость не более $n+1$, а $p+q\le n-t$. Поэтому
\[
T(n)\le 1+2t(n-t)+t(t+1)+(n+1)
      =2tn-t^2+t+n+2.
\]
При $t\le \lfloor n/3\rfloor-1$ эта граница меньше значений,
указанных ниже. Следовательно, достаточно сравнить случаи с
$k=\lfloor n/3\rfloor$ полными итерациями.

Если $n=3k$, все символы должны входить в блоки длины $k$:
$p=q=r=k$, и входное слово равно $a^k b^k c^k$.
Если $n=3k+1$, дополнительный символ можно добавить к одному из блоков
или оставить после $a^k b^k c^k$. Подстановка в точные формулы даёт
\[
\begin{array}{c|c}
(p,q,r;\,|v|) & T(n)\\ \hline
(k+1,k,k;0) & 5k^2+5k+3\\
(k,k+1,k;0) & 5k^2+5k+2\\
(k,k,k+1;0) & 5k^2+4k+2\\
(k,k,k;1) & 5k^2+4k+2
\end{array}
\]
Следовательно, добавочный символ выгоднее всего поместить в блок $a$.

Если $n=3k+2$, распределение двух дополнительных символов даёт
следующие наибольшие конкурирующие значения:
\[
\begin{array}{c|c}
(p,q,r;\,|v|) & T(n)\\ \hline
(k+1,k+1,k;0) & 5k^2+8k+4\\
(k+2,k,k;0) & 5k^2+7k+4\\
(k,k+2,k;0) & 5k^2+7k+2\\
(k+1,k,k+1;0) & 5k^2+5k+3\\
(k,k+1,k+1;0) & 5k^2+5k+2\\
(k,k,k+2;0) & 5k^2+4k+2\\
(k+1,k,k;1) & 5k^2+5k+3\\
(k,k+1,k;1) & 5k^2+5k+2\\
(k,k,k+1;1) & 5k^2+4k+2\\
(k,k,k;2) & 5k^2+4k+2
\end{array}
\]
Здесь $|v|$ --- длина суффикса после блока $c^r$; если он непуст,
его первый символ равен $a$ или $b$. При $k\ge1$ максимум достигается
на $p=q=k+1$, $r=k$.

Итак, точная сложность в худшем случае равна:
\[
\begin{array}{c|c|c}
 n & \text{Слово максимальной сложности} & T_{\max}(n)\\ \hline
3k,\ k\geq1 & a^k b^k c^k & 5k^2+7k+3\\
3k+1,\ k\geq1 & a^{k+1}b^kc^k & 5k^2+5k+3\\
3k+2,\ k\geq1 & a^{k+1}b^{k+1}c^k & 5k^2+8k+4
\end{array}
\]
Для малых длин: $T_{\max}(1)=4$ (слова $b$ и $c$),
$T_{\max}(2)=4$ (например, $aa$ или $ab$).

\subsection*{Точная сложность в среднем случае}

Предположим, что все слова из $\{a,b,c\}^n$ равновероятны.
Пусть $S_n$ --- сумма числа переходов по всем таким словам. Тогда
\[
T_{\mathrm{avg}}(n)=\frac{S_n}{3^n}.
\]
Для пустого слова $S_0=2$, а при $n=1$ времена для слов $a,b,c$
равны соответственно $3,4,4$, поэтому
\[
T_{\mathrm{avg}}(0)=2,\qquad T_{\mathrm{avg}}(1)=\frac{11}{3}.
\]

Для записи точной формулы введём сумму по трём положительным блокам.
Для целого $m$ положим
\[
H_m^{\mathrm{off}}=
\sum_{\substack{p,q,r\geq1\\p+q+r=m}}
T_{\mathrm{off}}(p,q,r),\qquad
H_m^{\mathrm{end}}=
\sum_{\substack{p,q,r\geq1\\p+q+r=m}}
T_{\mathrm{end}}(p,q,r).
\]
При $m<3$ обе суммы пусты и считаются равными нулю.
При $n\geq2$ точная сумма времён равна
\[
\begin{aligned}
S_n={}&2\cdot3^n
 +\sum_{p=1}^{n-1}3^{n-p-1}(p+2)+(n+2)\\
&+\sum_{m=2}^{n-1}(m-1)3^{n-m-1}(m+2)
 +(n-1)(n+2)\\
&+2\sum_{m=3}^{n-1}3^{n-m-1}H_m^{\mathrm{off}}
 +H_n^{\mathrm{end}}.
\end{aligned}
\]
Пустые суммы считаются равными нулю. Формула получена разбиением всех
слов на непересекающиеся классы:
\begin{itemize}
\item слова, начинающиеся с $b$ или $c$; при $n\geq2$ каждое из них
      завершается за $3$ перехода;
\item слова вида $a^p c v$ и слово $a^n$; их время равно $p+2$
      и $n+2$ соответственно;
\item слова вида $a^p b^q a v$ и слова $a^p b^q$;
      их время равно $p+q+2$;
\item слова вида $a^p b^q c^r v$ с $p,q,r\geq1$, где либо
      $v=\varepsilon$, либо $v$ начинается с $a$ или $b$.
      Их время задаётся функциями $T_{\mathrm{end}}$ и
      $T_{\mathrm{off}}$.
\end{itemize}
В третьем классе при $p+q=m<n$ число продолжений равно
$3^{n-m-1}$, а при $m=n$ имеется одно слово на каждое разбиение
$m=p+q$. В последнем классе число непустых продолжений для фиксированных
$p,q,r$ равно $2\cdot3^{n-p-q-r-1}$; это объясняет коэффициенты
в точной формуле.

Следовательно, приведённые выражения задают точное число переходов
в худшем случае и точное среднее число переходов для каждого $n$;
это не только асимптотические оценки.

\end{document}